%!TEX root = ../thesis.tex
% Requires the bibliography and cross-reference packages already loaded by the university template.

\chapter{Discussion and Limitations}
\label{chap:discussion}

\ifpdf
    \graphicspath{{Chapter7/Figs/Raster/}{Chapter7/Figs/PDF/}{Chapter7/Figs/}{Chapter7/figures/}}
\else
    \graphicspath{{Chapter7/Figs/Vector/}{Chapter7/Figs/}{Chapter7/figures/}}
\fi

Chapter~\ref{chap:experimental-results} showed that generated Bangla cinema responses can be made more consistent with a requested engagement-specificity level, while also exposing a gap between successful control and faithful audience modelling. This chapter interprets that gap and considers the limits of the evidence. Contributions and future work are reserved for the final chapter.

\section{Main Finding}
\label{sec:ch7-main-finding}

The evidence shows that conditioning and revision make short Bangla cinema responses more consistent with a requested engagement-specificity level than zero-shot prompting. Verifier-B provides consistent coverage of the complete frozen generation surface, and native Bangla readers usually identify the same distinction on a smaller blinded subset. Neither result demonstrates that the system reproduces a real audience. The generated responses were not matched against reactions to the same films, and the two levels do not represent demographic or psychological groups. The result is therefore about control over a textual distinction, not prediction of viewers, sentiment, popularity, or commercial performance. This boundary is especially important because synthetic-audience research is vulnerable to prompt sensitivity, hallucination and anthropomorphic interpretation~\citep{b1}. No single mechanism accounts for the improvement: static examples, retrieval, blind resampling, self-critique and verifier-guided revision all improve on zero-shot under the planned comparisons. The neuro-symbolic condition has the largest observed difference from zero-shot, but the experiment does not rank the active systems. Its direct comparison with the neural-only loop was post-hoc and remains exploratory.

\subsection{Engagement Specificity as a Continuum}
\label{sec:ch7-continuum}

The early corpus analysis changed the meaning of the target variable. The Region-A partition was reproducible, yet the geometric evidence did not support two separated clusters: the silhouette was weak, the gap statistic selected no value of $K$, and HDBSCAN treated the points as noise. Human comparative judgment showed that readers could nevertheless recognize the contrast and its direction. This combination supports a usable cut through a continuum, not two naturally occurring audience types.

The framework can therefore request a broader or more specific response without assigning a person to a fixed category. It controls a property of the generated text; it does not infer personas within Bangla-speaking audiences.

Length remains part of the construct problem. The corpus, development outputs and final generations all show that engagement specificity and response length are related, although the direction is not stable across every setting. The length-matched analysis retains a Level-0 deficit, so length alone does not explain the result. At the same time, matching keeps only a selective fraction of the outputs and cannot establish length-neutral control. The model appears to manipulate a bundle of cues---detail, reference to plot events, evaluative specificity and length---rather than an isolated semantic property.

\section{What the Control Mechanisms Add}
\label{sec:ch7-control-mechanisms}

The ablations separate four functions that would otherwise be hidden inside one system score:

\begin{itemize}
    \setlength{\itemsep}{2pt}
    \setlength{\parskip}{0pt}
    \setlength{\parsep}{0pt}
    \item \textbf{Examples and retrieval.} Static demonstrations and R1 retrieval show the two levels in Bangla. They improve the prompt but do not decide when a draft has reached its target.
    \item \textbf{Neural verification.} Verifier-A supplies the stopping decision. Revision costs additional calls, and later attempts contain only earlier failures, so their trajectory describes difficult survivors rather than a fixed population.
    \item \textbf{Symbolic diagnosis.} Symbolic acceptance is weak and inefficient, particularly at Level 1. Symbolic feedback is more useful for naming a problem while the neural verifier retains the stopping decision. The Level-0 increment over neural-only is exploratory, not evidence of general hybrid superiority.
    \item \textbf{Alternative critics.} Self-critique improves control at a fixed three-call cost. The hosted judge is a comparator, not an independent outcome authority. Evaluator replacement, self-preference and uneven multilingual agreement reinforce this restriction~\citep{b22,b23,b24,b25}.
\end{itemize}

\subsection{Why Verifier Separation Matters}
\label{sec:ch7-verifier-separation}

Verifier separation becomes most informative during revision. In the neural-gated loops, drafts selected to satisfy Verifier-A receive progressively less comparable support from Verifier-B. Because both scores are available on the same continuing cases, the divergence cannot be attributed to a change in plot or requested level. Improvement against the in-loop proxy and improvement under a held-out instrument are not interchangeable.

This pattern is consistent with Goodhart-style pressure, but the evidence does not identify reward hacking in every revised response. Later attempts are failure-selected, Verifier-B has its own errors, and its calibration improvement was not established. Repeated optimization nevertheless changes the relationship between the two instruments in the direction expected when a proxy becomes increasingly exploitable. Evaluator stress-testing work similarly argues that an optimized score needs independent or invariant checks before it can stand for the intended construct~\citep{b7}.

Keeping Verifier-B outside generation therefore does more than prevent direct leakage. It makes disagreement observable. Had the same verifier controlled revision and supplied the final outcome, the experiment could report rising scores without showing whether the revisions generalized to another trained instrument. The separation does not make Verifier-B ground truth, but it turns proxy dependence into an empirical question.

\section{Human and Automatic Evidence}
\label{sec:ch7-human-automatic}

The blinded study tests the meaning of the output labels rather than providing another condition leaderboard. Native Bangla readers usually recover the requested level with strong panel agreement, but five items per condition-level cell are too few for reliable system ranking.

The same-item comparison reveals an important asymmetry. Human majority labels perform similarly at both levels, whereas both automatic verifiers are weaker at Level 0. This makes an instrument-specific Level-0 problem plausible. It does not prove that the human panel is the definitive oracle: the comparison contains only 100 items, and within-level agreement statistics are affected by fixed requested-label marginals~\citep{b36,b37}. Still, the discrepancy warns against describing Verifier-B accuracy as human quality.

Annotators selected the requested level; they did not rate fluency, plot faithfulness, sentiment, usefulness, offensiveness or resemblance to real audience responses. Human validation therefore supports controllability of the stated distinction, not overall response quality.

\section{Practical Implications}
\label{sec:ch7-practical-implications}

The evidence supports three practical decisions:

\begin{itemize}
    \setlength{\itemsep}{2pt}
    \setlength{\parskip}{0pt}
    \setlength{\parsep}{0pt}
    \item \textbf{Appropriate use:} an inspectable pre-writing aid for comparing both levels, reviewing retrieved examples and revision status, and preparing questions for direct audience research.
    \item \textbf{Unsupported use:} autonomous marketing, casting, segmentation or commercial forecasting. The study learns no film-level response distribution and links no generated statement to a demographic group.
    \item \textbf{Deployment choice:} static prompting, retrieval or blind resampling may suit inexpensive variety; neural-guided loops add stopping and correction traces. The symbolic-only loop is not supported as a default, and cost should not be hidden inside one quality score.
\end{itemize}

\section{Limitations}
\label{sec:ch7-limitations}

\subsection{Construct and Measurement}
\label{sec:ch7-construct-measurement}

\begin{itemize}
    \setlength{\itemsep}{2pt}
    \setlength{\parskip}{0pt}
    \setlength{\parsep}{0pt}
    \item The axis comes from one Bangla review corpus with unresolved source history; region, sentiment and length complicate its geometric origin. Human validation supports an operational label, not a natural audience taxonomy.
    \item Verifier-B is a fixed instrument, not a calibrated probability of human acceptance. Its Level-0 mismatch and calibration null remain unresolved.
    \item The realism measures capture different properties. LaBSE-feature MAUVE is small-sample sensitivity evidence~\citep{b55}, and sentiment divergence remains unmeasured because no independent registered scorer was available.
\end{itemize}

\subsection{Internal and Statistical Validity}
\label{sec:ch7-internal-statistical}

\begin{itemize}
    \setlength{\itemsep}{2pt}
    \setlength{\parskip}{0pt}
    \setlength{\parsep}{0pt}
    \item Verifier-A and Verifier-B use disjoint data and different model families, but learn the same operational label; their errors need not be independent. The hosted judge is also a same-family treatment, not an external authority.
    \item Later attempts are failure-selected, and length matching conditions on a property changed by generation. Neither analysis supports an unrestricted causal reading.
    \item The primary tests use 540 paired cases, paired bootstrap, McNemar testing and Benjamini--Hochberg correction~\citep{b49,b50}. The three seeds are sensitivity blocks, not independent replications, and the direct hybrid contrast remains outside the registered comparison family.
\end{itemize}

\subsection{External Validity, Human Evaluation, and Ethics}
\label{sec:ch7-external-human-ethical}

\begin{itemize}
    \setlength{\itemsep}{2pt}
    \setlength{\parskip}{0pt}
    \setlength{\parsep}{0pt}
    \item The corpus has no film identifier. The 90 held-out plots separate evaluation from development but do not establish generalization to other Bangla registers, communities, domains or generator families.
    \item The three adult native-Bangla evaluators were university batchmates known to the researcher. Blinding and coded identities do not remove the limits of a convenience sample or possible social pressure. No institutional approval or exemption is claimed, and absent fields were not reconstructed~\citep{b32}.
    \item Outputs may contain stereotypes, offensive language or unsupported plot details. The interface's source-support diagnostic was not part of the frozen experiment or a registered human faithfulness audit; it neither establishes a hallucination rate nor certifies safety.
\end{itemize}

\section{Chapter Summary}
\label{sec:ch7-summary}

The study shows that engagement specificity can be varied in short Bangla cinema responses, while the verifier and human evidence explains why this result is not audience simulation. Chapter~\ref{chap:conclusion} closes the argument and identifies the evidence still needed beyond this corpus.
